
    \section {资产概述}

.
    \begin{center}
        \includegraphics[scale=0.2] {gold.jpg}
    \end{center}

    会计学上的资产（Asset），指一企业透过交易或非交易事项所获得之经济资源，能以货币衡量，并预期未来能提供效益者。

\subsection {资产的属性}

    依据资产的一般定义，它具有下述特性：
    \begin{itemize}
        \item  该事物存在未来的经济价值。“未来的经济价值”是指，特定企业（或组织）若拥有该事物（资产），则可以利用其来产生（增加）企业未来的现金流入，或利用其替代（减少）企业未来发生的现金流出。
        \item  该事物产权被特定企业控制。
        \item  若某事物属于资产（存在未来的经济价值），但非某企业可控制，则某企业无法利用其经济价值，故该事物对于某企业不属于资产。
        \item  若该事物属于公共财产，虽然该事物对特定企业可能存在未来的经济价值，但公共财产非被特定企业所控制，故公共财产皆不为特定企业的资产。
    \end{itemize}

    资产是财务会计中争议最大的概念之一。历来的会计学家都试图对资产给出满意的界定，但到目前为止，一个权威的、被学术界和实务界所共同认可的定义，尚未出现。

    资产的经济属性即能够为企业提供未来经济利益，这也是资产的本质所在。也就是说，不管是有形的还是无形的，要成为资产，必须具备能产生经济利益的能力，这是资产的第一要义。

    资产的法律属性即必须是为企业所控制，也就是说，资产所产生的经济利益能可靠地流入本企业，为本企业提供服务能力，而不论企业是否对它拥有所有权，这是资产的第二要义。　　

\subsection {资产的分类}

    财务会计中的资产是指预计有助于生产未来现金流入或减少未来现金流出的经济资源，一般分为：

    \begin{itemize}
        \item  依据是否有具体形体分类
            \begin{itemize}
                \item  有形资产
                \item  无形资产
            \end{itemize}

        \item  依据耐用期限的长短分类
            \begin{itemize}
                \item  长期使用资产（固定资产）
                \item  流动资产
            \end{itemize}

        \item  依据资产所有权分类
            \begin{itemize}
                \item  贷方权益（负债）
                \item  业主权益
            \end{itemize}
    \end{itemize}

\subsection {资产定义的几种观点}

    未消逝成本观：未消逝成本观是对资产性质的早期描述。美国著名会计学家佩顿和利特尔顿在《公司会计准则导论》(an introduction to corporate standard)(1940)中明确提出了未消逝成本观。他们认为：“……成本可以分为两部分，其中已经消耗的成本为费用，未耗用的成本为资产……”也就是说，他们认为资产是营业或生产要素获得以后尚未达到营业成本和费用阶段的那部分余额，是成本中未消逝的那部分余额。显然，这种观点同历史成本会计模式是密不可分的。它着重从会计计量的角度来定义资产，强调了资产取得与生产耗费之间的联系。

    借方余额观：资产定义的借方余额观是由美国会计师协会(美国执业会计师协会的前身)所属的会计名词委员会在其颁布的第1号《会计名词》(1953)中提出来的。该公告认为：“资产是由借方余额所体现的某种东西。这一借方余额是按照公认会计原则或规则从结平的各账户中结转过来的，前提是这一借方余额不是负值。作为资产，它代表的或者是一种财产权利，或者是所取得的价值，有的则是为取得财产权利或为将来取得财产而发生的费用支出”。这一认识的基本特征是将资产视为借方余额的体现物。据此，不仅借方余额所体现的应收账款、存货、设备、厂房等要确认为资产，而且由借方余额所体现的递延费用等项目也可以确认为资产。显然，这种观点只是从会计结账技术的角度来理解资产，很难说是在描述资产的性质。

    经济资源观：经济资源观是关于资产定义的颇具影响的一种观点。1957年，美国会计学会发表的《公司财务报表所依恃的会计和报表准则》中明确指出：“资产是一个特定会计主体从事经营所需的经济资源，是可以用于或有益于未来经营的服务潜能总量”。对资产的这一认识，第一次明确地将资产与经济资源相联系，虽然它并未正面提到无形资产的内容，但这一定义至少可能将无形资产包纳其中。另外，它也明确了资产与特定会计主体之间的关系，即特定会计主体能够借助资产业从事未来经营。1970年，美国会计准则委员会在其发布的第4号公告中提出了一个资产定义：“资产是按照公认会计原则确认和计量的企业经济资源，资产也包括某些虽不是资源但按照公认会计原则确认和计量的递延费用”。这一定义虽然明确指出资产的实质是经济资源，接受了“经济资源”这一新认识，但它却认为，经济资源应否视为资产，取决于公认会计原则的确认和计量标准，这就把资产的实质与资产的确认和计量之间的主从关系颠倒了。

    未来经济利益观：目前比较流行的资产定义体现了未来经济利益观的观点。1962年，穆尼茨(moonitaz)与斯普劳斯(r.t.sprouse)在《会计研究论丛》第3号——《企业普遍适用的会计准则》这一文献中明确提出：“资产是预期的未来经济利益，这种经济利益已经由企业通过现在或过去的交易获得。”现在的美国财务会计准则委员会(fasb)在《财务会计概念公告》第6号(sfac no.6)中提出：“资产是可能的未来经济利益，它是特定个体从已经发生的交易或事项中所取得或加以控制的。”

    未来经济利益观认为，资产的本质在于它蕴藏着未来的经济利益。因此，对资产的确认或判断不能看它的取得是否支付了代价，而要看它是否蕴藏着未来的经济利益。在现实中，虽然成本是资产取得的重要证据之一，而且成本还是资产计量的重要属性，但是，成本的发生并不一定导致未来的经济利益，而未来经济利益的增加也并不必然会发生成本，例如，业主投资、接受捐赠等。所以，未消逝成本观将未耗用的成本看成是资产，视资产为成本的组成部分，是不切实际的。而经济资源观强调资产的经济资源属性，把一些不是经济资源但有助于实现未来经济利益的或减少未来经济损失的项目如某些备抵项目排斥除在资产之外。未来经济利益观则将这些项目合乎情理地包括在资产之中。因此，我们说未来经济利益相对于其他观点来说更加全面、合理。
 % 资产概述
    \section {财务预算}

\subsection {新阶段，实行财务预算管理的重要性}

    就我国目前企业经营管理的实际成效调查分析来看，部分企业中的财务管理仍处于低层次管理阶段，甚至在一些企业中仍未使用实际性、可持续操作的预算管理办法。据调查，在大型企业经营管理中财务预算管理越来越发挥着重要的作用。新阶段，随着我国企业种类、企业数目的扩大，财务预算管理的普及势在必行。

  \subsubsection { (一)有利于提升企业的整体竞争力}

    财务预算管理是企业管理的核心，而预算就是以货币表示的预算结果，既是决策的具体化，又是控制生产经营的活动的依据。它把企业经营管理的计划与控制结合了起来：一方面为做出这份方案提供可靠而充实的资金的计划;另一方面对该方案进行监督与考核，权衡利弊。

   \subsubsection {(二)有利于强化企业资金管理，其本质就是获得资源利用的最大化}

    财务预算管理是实现对产品流和资本流的管理控制，其目的在于用最低的投入获得最高额的收益。通常，就是使得日常资金预算达到资金运用的最高效益。在此基础上，营业收入预算是整个财务预算管理的中枢环节，营业收入关系到整个预算方案的合理性和可持续性。

    \subsubsection {(三)有利于评价各部门各单位工作绩效}

    财务预算是企业在预测和决策的基础上，围绕着战略计划，对年度预算内各种经济资源和经济行为的合理预计、测算并进行的财务和监督的活动。其侧重点更在于对各部门年度工作成绩总结的一种考核，对其部门能力的一种鉴定。

    \subsubsection {(四)推行财务预算管理，促进企业管理创新是企业发展的必然选择}

    现代企业经营理念更注重的是“产权明晰、权责明确、管理科学”，企业财务管理内部功能、外延及其地位发生了深刻的变化，强化财务预算管理已成为现代管理企业在激烈的市场竞争中得以生存和发展的一大法宝，也是现代企业管理制度得以实施的重要保障。企业应当积极按照国家相关规定，组织好财务预算工作，培养相关人才，明确责权限，完善编制程序，强化监督，从内因解决问题。

\subsection {财务预算管理在企业管理中的现状分析}

   \subsubsection { (一)对财务预算管理观念淡薄，财务管理执行力弱}

    就目前来看，不少电缆加工企业管理者把企业生存发展的重中之重放在产量上面，追逐订单量，也就是交易量，过于看重营业收入而不考虑各个方面的隐性成本，比如说对环境卫生影响、质量不过关召回等所另外担负的成本。此外，从业职工培训不过关，缺乏相关理论知识的指导，使得财务预算部门的执行力大打折扣。

    \subsubsection {(二)收入预算、支出预算达不到统一}

    在实际生产过程中，很多企业很难以做到一方面挖掘足够的收入，清理欠款并合理占用资金;另一方面紧紧抓住资金管理的关键点，全面实施财政收支预算管理。预算已经确立就相当于具有了权威性，换句话说就是“法律效力”，正确的权威性才能引领企业走向成功，而当今不合理的预算依然存在，对企业的生存发展构成一定的威胁。

    \subsubsection {(三)财务预算工作已完成上级部门下达的任务为目标，预算执行往往超标}

    据了解，很多企业在财务预算方面对各部门采取比较传统的承包式的模式，最终预算执行往往超标，几经调整也无济于事，特别在年终考核，实际日数与预算期限大有出入，收入不能确保，成本费用得不到很好地控制，尤其是利润增长缓慢出现明显停滞现象。
 % 财务预算
        现金流量表作为企业财务报告的重要组成部分，是反映企业“血液循环”——现金及其等价物的增减及变化后的状况的动态报表，具有鲜明的综合性、可比性。但是原始的财务报表数据资料只有科学的运用专门方法，加以疏理、分析，才能达到揭示企业财务状况变动的原因、后果及趋势，为报表使用者预、决策提供正确的依据的目的。就现金流量表而言，可运用结构分析法、比率分析法和趋势分析法三种分析方法。

\subsection {现金流量表的主要内容}

    \begin{enumerate.zh}
        \item  概况 分析现金流量采用的最基本的方法是编制百分数现金流量表。即分别计算出各项现金流入、流出及净值占现金总流入、流出及净值的比重。通过结构百分比分析，可反映出企业现金流入（流出）的渠道流出有几条，要增加现金流入主要靠什么；流出的渠道分别占有多大比例，通过分析了解企业财务状况的形成、变动及原因等。一般来讲，经营活动现金流占总现金流比重大的企业，经营状况较好，财务风险较低，现金流入结构也较为合理。特别当现金净流量为正时，判断企业现金流入是否强劲，要注意现金净流量是由经营活动产生的还是融资活动产生的，从而深入探究经营活动产生的现金流的源泉是否稳定、可靠。

        \item  经营活动产生的数据流分析 进行此项分析时，可对经营活动产生的数据流细分项进行结构分析，同利润表中的主营业务收入和其他经营利润相结合，若两者相差不大，说明企业会计帐面上的收入数值非常有效、及时地转成了实际现金流入，应收账款的管理也比较有效（可参考应收账款周转率分析），未来为保证其经营活动产生的现金源源流入提供了有力保障；反之，则应加强经营收入的实现与变现，确保资金的有效运转，否则，即使账面资产较大，也存在潜亏的可能。

        \item  投资活动产生的现金流分析 进行此项分析时，应充分考虑企业预算、投资计划，可对应资产负债表中的长期投资及历年投资收益等相关情况进行分析，判断是否存在潜在风险。例如，企业投资活动现金流出是用于购建资产且自有资金足以支付的，可分析应是按计划扩大生产规模或改进生产工艺，以降低经营成本，预测会为企业发展奠定良好的基础。

        \item  融资活动产生的现金流分析 此项分析可反映出企业融资来源与用途及结构比率，一般与现金流量表概况分析相结合，考虑现金流入、流出的总体趋势，对举借债务的结合比率分析法进一步分析偿债能力。
    \end{enumerate.zh}

\subsection {现金流量表的分析方法}

    现金流量表比率分析是以经营活动现金净流量与资产负债表等财务报表中的相关指标进行对比分析，全面揭示企业的经营水平，测定企业的偿债能力，反映企业的支付能力。大致可分为盈利质量分析、偿债能力分析和支付能力分析三部分。具体分析如下：

    \begin{enumerate.zh}
    \item  盈利质量分析。

    盈利质量分析是指企业根据经营活动现金净流量与净利润、资本支出等之间的关系，揭示企业保持现有经营水平，创造未来盈利能力的一种分析方法。包含以下指标：

        \begin{itemize}
            \item  盈利现金比率：=经营现金净流量/净利润。一般来讲，该比率越大，企业盈利质量越强。若该比率<1，说明企业本期净利润中存在未实现的现金收入，即使盈利，也可能发生现金短缺，严重时亦可导致亏损。分析时可参考应收账款，若本期增幅较大，需及时改进相关政策，确保应收款按时回收，消除潜在风险。

            \item  再投资比率：=经营现金净流量/资本性支出，反映企业当期经营净现金流量是否足以支付资本性支出（固定资产的投资）所需现金。一般来讲，该比率越高，说明企业扩大生产规模、创造未来现金流量或利润的能力就越强，若<1，说明企业资本性支出包含融资以弥补支出的不足。
        \end{itemize}

    \item  偿债能力分析。

    企业经常通过举债来弥补自有资金的不足，但最终用于偿债的最直接的资产是现金，因此，用现金流量和债务比较可以更好的反映企业的偿债能力。

        \begin{itemize}
            \item  现金流动负债比：=现金净流量/流动负债总额，若此指标偏低，反映企业依靠现金偿还债务的压力较大；偏高，则说明企业能轻松的依靠现金偿债。

            \item  现金债务总额比：=现金净流量/债务总额，该指标反映当年现金净流量负荷总债务的能力，可衡量当年现金净流量对全部债务偿还的满足程度。分析时应与债务平均偿还期相结合，若平均偿还期越短，则比率越高越好，反之则此比率要求越低。企业债权人可凭此指标衡量债务偿还的安全程度。

            \item  现金到期债务比：=现金净流量/本期到期债务。这里，本期到期债务是指即将到期的长期债务和应付票据。对这个指标进行考察，可根据其大小直接判断公司的即期偿债能力。
        \end{itemize}

    \item  支付能力分析。

    支付能力的分析，主要是通过企业当期取得的现金收入特别是其中的经营活动现金收入，同其各种开支来进行分析和比较的。

    将企业本期经营活动现金收入同本期偿还的债务、发生支出进行对比后，其余额即为可用于投资及分配的现金。在不考虑筹资活动的情况下，它们的关系是：可用于投资、分配股利（利润）的现金=本期经营活动的现金收入+投资活动取得的现金收入-偿还债务的现金支出-经营活动的各项开支。

    一般来讲，若企业本期可用于投资、分配股利（利润）的现金>0，说明企业当期经营活动现金收入加上投资活动现金收入足以支付本期债务及日常活动支出，且尚有结余用于再投资或利润分配；反之指标<0，说明尚需通过筹资来弥补支出的不足。

    \end{enumerate.zh}

\subsection {从现金流量看企业的发展趋势}

    现金流量表趋势分析主要是通过观察连续报告期（至少为2年，比较期越长，越能客观反映情况及趋势）的现金流量表，对报表中的全部或部分重要项目进行对比，比较分析各期指标的增减变化，并在此基础上判断其发展趋势，对未来做出预测的一种方法。趋势分析注重可比性，具体问题具体分析。例如，正常经营的同一企业在不同时期如果采用不同的财务政策，其现金流量表的变化并不能完全说明其财务状况的变化趋势；筹建期或投产期的企业，经营活动现金流量较少，筹资活动的现金流量较多，成熟期一般则相反。因此，正确计算运用趋势百分比，可使报表使用者了解有关项目变动的基本趋势及其变动原因，在此基础上预测企业未来的财务状况，为其决策提供可靠的依据。

\textbf {小结}

    综上论述，我们可以看出，现金流量表在分析企业财务状况时，确实是一个不可多得的工具。在实际操作中，注意现金流量表分析与资产负债表和利润表等财务报表分析相结合，可以更清晰、全面地了解企业的财务状况及发展趋势，了解其与同行的差距，及时发现问题，正确评价企业当前、未来的偿债能力、支付能力，以及企业当前和前期所取得的利润的质量，科学的预测企业未来财务状况，为报表使用者做出决策的提供正确的依据。
 % 资金的现金流量表
    \section {资金的现金流管理}

    企业是以盈利为目的的，是人尽皆知的。当前不乏有一些企业刻意的追求高收益、高利润。因此往往会有这样一种错误的思想，认为企业利润显示的数值高就是经营有成效的表现。从而很大一定程度上忽略了利润中所应该体现出来的流动性。作为企业的资金管理者应当要能够充分、正确地界定现金与利润之间的差异。利润并不代表企业自身有充裕的流动资金。流动资金也就是我们俗称的现金流，对于企业的健康发展有着重要的作用。

    现金流（Cash Flow）是指一段时间内企业现金流入和流出的数量。企业在销售商品、提供劳务、或是出售固定资产、向银行借款的时候都会取得现金，形成现金的流入。而企业为了生存、发展、扩大需要购买原材料、支付工资、构建固定资产、对外投资、偿还债务等，这些活动都会导致企业现金流的流出。如果企业手头上没有足够的现金流来面对这些业务的支出，其结果是可想而知的。从企业整体发展来看，现金流比利润更为重要，它贯穿于企业的每个环节。在现实生活中我们可以看到，有些企业虽然帐面盈利颇丰，却因为现金流量不充沛而倒闭；有的企业虽然长期处于亏损当中，但其却可以依赖着自身拥有的现金流得以长期生存。企业的持续性发展经营，靠的不是高利润而是良好、充足的现金流。

    传统意义上的现金管理主要是涉及企业资金的流入流出。然而广义上的现金管理，其所涉及的范围就要广的多，通常包括企业帐户及交易管理、流动性管理、投资管理、融资管理额风险管理等。企业现金管理主要可以从规划现金流、控制现金流出发。

    规划现金流主要是通过运用现金预算的手段，并结合企业以往的经验，来确定一个合理的现金预算额度和最佳现金持有量。如果企业能够精确的预测现金流，就可以保证充足的流动性。企业的现金流预测可以根据时间的长短分为短期、中期和长期预测。通常期限越长，预测的准确性就越差。到底选择何种现金流的预测方式就要纵观企业的整体的发展战略和实际要求。同时企业的现金流预测还可以现金的流入和流出两方面的出发，来推断一个合理的现金存量。

    控制现金流量是对企业现金流的内部控制。控制企业的现金流是在正确规划的基础上展开的，主要包括企业现金流的集中控制、收付款的控制等。现金的集中管理将更有利于企业资金管理者了解企业资金的整体情况，在更广的范围内迅速而有效地控制好这部分现金流，从而使这些现金的保存和运用达到最佳状态。
 % 资金的现金流管理
        “现金为王”的理念在当前企业资金管理中倍受关注。随着企业集团全球化业务的不断拓展，企业对于内部资金的集中化管理和流动性管理的需求变得日益迫切。企业资金管理正逐渐从成本中心向利润中心转换。管理好企业的现金流已经不单单拘泥于管理好库存现金的层面，更多的指的是营运资金。如何强化企业的现金流管理，提升企业价值成为了越来越多资金管理人的工作重点。近年来，随着我国监管体系的不断放宽，一种基于委托贷款模式的创新型资金管理技术—“现金池”正逐渐兴起，成为了集团企业实现资产流动性、集中化管理的首选工具。

    什么是企业的现金池？

    现金池是建立在委托贷款模式下的集团企业内部的资金融通。在现金池结构中有一个母帐户和若干个子帐户。以公司总部的名义设立集团现金池母帐户，通过子帐户向母帐户以委托贷款的方式，每日定时将子帐户的资金余额上划到母帐户中。资金上划后，子公司账户上保持零余额或目标余额。这个限额的设定通常是由企业根据自身资金管理的需求和现金存量的额度与银行协商确定。在现金池结构中，子账户资金上划一般是在固定的时点，一天一次。但是企业也可以根据实际需求向银行要求一天多次的子账户余额上划。

    日间，若子公司对外付款时账户余额不足，银行可以根据以其上存母账户的资金头寸额度为限进行透支支付。一般而言，为企业提供现金池业务的银行都会为其子账户提供日间透支服务，以确保子账户日间对外的正常支付。当然，银行为成员单位提供日间透支服务，是基于企业集团自身良好的信誉和相应的资产担保，必要时企业还需事先和银行签订相关的法律文书。银行规定，集团企业为其下属子账户提供担保的必须要大于或者等于其所有子账户日间透支额度的总和。子账户的日间透支会由现金池的母账户在日终自动向子账户划拨相应的款项来结清欠款。企业集团还可以根据每个子账户的资金需求情况合理分配他们的日间透支额度。只要子帐户的日间透支额度不超过企业集团担保的额度，银行是不会就这一透支来征收利息的。但是在实际操作中，往往会有突发情况的发生，子账户对外支付金额可能要超过其日间可透支额度。这就需要企业集团向银行发出人工拨付的请求，由银行将款项拨付到需要付款的子账户。这样企业集团既能及时满足子账户应急的需要，也不会使子账户支付需求超出额度，而承担额外的利息费用。这种临时性的拨付一般是由企业人员通过网上操作完成的，计算机网络的发达，系统化，足以使得这类款项实时到帐。

    企业现金池业务的基本流程图

    以银行和财务公司为委托贷款中介的现金池

    企业现金池可以以银行，也可以以企业集团财务公司为委托贷款中介进行开设。现金池业务中涉及到的事项主要有子帐户余额的上划、子帐户的日间透支、主动拨付、母子帐户之间的委托贷款，以及子帐户向集团总部的上存、下借分别计息等。如下图所示：
    　　

    以银行为委托贷款中介

    在此模式下的现金池，集团各成员企业均需在现金池服务提供的银行开立结算账户、委托存款账户和委托贷款账户，其中结算账户为法人透支账户，具体透支额度可由银行对各家成员企业逐个单独授信，也可由集团核心企业统一申请和担保的集团授信并由核心企业将授信额度分配给各成员企业，授信提款可采用法人账户隔夜透支的方式，也可采用发放一笔流动资金贷款的方式。

    集团企业和法人关联公司的账户超过目标余额的部分，首先归还最早借入的委托贷款，其次划拨到委托存款账户。如果出现导致现金池中的汇总委托存款余额小于汇总委托贷款余额的借入委托贷款请求，委托存款余额小于委托贷款余额的部分形成资金需求企业的账户透支。

    在该种现金池中，集团财务公司可作为一个成员企业的角色出现，并代表集团各成员企业与银行沟通。当委托存款虚拟汇总账户余额大于委托贷款汇总余额的部分或存差，该差额部分可用于隔夜投资。当经常性出现委托存款虚拟汇总账户余额大于委托贷款余额，其差额中部分稳定的余额，可由财务公司作为牵头企业代表集团成员投资于定期存款、通知存款、协议存款，或投资于国债市场、货币市场、基金市场，或向其他企业发放委托贷款。

    以财务公司为委托贷款中介

    其模式与以银行为委托贷款中介的现金池一致。由于财务公司支付中介的局限，财务公司一般需要选择一家银行作为结算代理，其成员企业也需要在银行开户与集团之外的其他企业进行结算。

    银行为财务公司提供现金池代理结算服务，此时银行并不是以委托贷款中介出现，而是以支付结算中介和或者现金池软件供应商的身份出现。由于委托贷款由财务公司处理，集团内的法人公司将富余资金划到在财务公司开立的委托存款专户账户，需要资金的企业在委托存款大于委托贷款的金额内借入委托贷款，银行按照财务公司和成员企业的指令划拨款项。如企业借入的委托贷款不足支付请求的部分，可由财务公司提供透支服务，也可通过银行或财务公司提供流动资金贷款解决。

    现金池业务给企业资金管理带来的积极作用

    企业通过建立现金池主要有三方面的好处：

    1. 利息成本的降低

    现金池实现了企业内部资金的融通，将外部融资变成了内部融资。在现金池中，不同帐户上的正负余额可以进行有效的相互抵消，资金充裕的帐户资金就能自动转移到那些资金缺乏的帐户上去，这样一来，既简化了融资手续，又降低了整个企业集团对于贷款的需求。而且内部融资的成本要远远低于银行的同期的贷款利率，这样也大大降低了融资成本和财务费用。但是为了避嫌涉及转移定价，内部融资的利率也应当遵守公平交易原则，采用的利率不宜过高也不宜过低。

    我们来举个例子，假设A，B，C同属一个集团公司，A日均盈余RMB400，000，B日均盈余RMB200，000，而C日均透支RMB300，000。我们假设现金池中的存款利率为2%，贷款利率为5%。很明显，集团通过设立现金池较之后就能够获得较多的利息收益。（我们暂且不计算相关税费，仅从利息收支方面出发）

    2. 获取最大的资金回报

    现金池将企业闲散的资金进行了回笼，充分提高了资金的使用效率。企业可以将集中后的资金进行更有效的投资活动，如投资于定期存款、通知存款、协议存款，或投资于国债市场、货币市场、基金市场，或向其他企业发放委托贷款等，从而为企业增加额外的投资收益。即使企业不进行投资活动，大额的存款也可以使企业从银行这里获得较高的协定存款利率。

    3.改善整个集团的资金管理水平

    在每日日终，各子帐户将余额上存母帐户，提高了整个集团资金的透明度。通过现金池，集团总部能够及时了解到各个子帐户的资金流量情况，便于企业司库们从整个集团层面管理和控制好流动性资金的需求，加强了内部控制的效力。
 % 现金池
        许多咨询公司认为，公司经营以人力资本为主，不需太多流动资金，更不需过高的资本投入。真是这样吗？

    现金流管理目前在生产型或商业流通型企业得到经营管理人员越来越多的重视，近日EVA概念在国内的推广者——思腾思特公司所面临着可能退出中国市场的尴尬局面，却不得不使许多公司引以为鉴，在现金流管理方面予以更多的重视。

    咨询公司很少提及现金流的问题，与思腾思特公司类似，典型的专业咨询公司，在现金流管理方面也和其他企业一样，存在问题。如果您的企业出现如下问题，那么，你的现金流已经亮起红灯。

    你的现金流能有效满足业务发展吗？

    ●资本金不足，无后续资金注入

    思腾思特公司，成立时启动资金仅一两百万美元，股东以财务性投资为主要动机，对于公司在运营中发生的现金流耗损没有通过进一步的注资来支持。许多咨询公司也都认为，公司经营以人力资本为主，不需太多流动资金，更不需过高的资本投入。孰不知，咨询公司运营流动资金的充足也是公司业务持续发展的基本条件，没有充足的现金流量，就似人体血脉供养不足，必将给公司持续经营带来问题。

    ●市场拓展并未带来足够的现金回流

    目前国内对咨询公司专业服务虽然十分认同，但对服务价值的认定却仍未充分考虑，许多咨询公司付出的专业劳动并不能获取充分的现金回报。思腾思特公司进入中国不久，就推出了“中国上市公司财富创造和毁灭排行榜”，并且以管理咨询、财务咨询等手段向中国不少国有大中型企业推广EVA概念及应用方法，市场增长潜力的确不小。然而由于处于市场拓展初期项目定价往往打有折扣，且客户来源于国内，咨询服务难度大，现金收益却来得低微，项目做了不少，年咨询业务收入却仅有1000余万元，扣除相关付现成本后，留存公司现金流量已很有限。

    ●缺乏其他有效的融资手段

    资本金较小，经营收益又不理想，且无足额的有效实物资产用于抵押，一般也就很难通过银行流动奖金贷款等其他融资手段解决资金缺口。

    现金支出是否难以控制？

    ●人力成本支出没有有效的规划

    许多咨询公司已经意识到人力资源对公司业务发展的重要性，但在人力成本支出方面却没有合理有效的规划，难以对员工形成有效的激励。员工对薪金收入不满，工作积极性不高，同时还会在公司财务控制不严的情况下设法增加自身各种费用开支，从而导致公司不能有效控制人力成本，现金支出不确定性增加，公司经营风险增大。

    ●项目运作成本没有有效的控制

    项目运作中，也没有有效的成本核算与支出控制措施，项目运作成本往往高出项目收入，导致在项目运作层面上的现金收支不平衡。

    ●公司财务控制体系并不健全

    咨询公司与生产型企业、商业流通性企业相比，规模要小许多。因而，公司往往忽视在财务控制体系方面的建设，现金预算制度、项目成本核算制度、费用报销制度等不甚健全，最终致使公司成本结束软化，现金开支难以预期与控制。现金流入的不足与现金支出的难以控制，会使咨询公司的现金趋于匮乏，最终将使公司正常的人力成本开支、项目成本开支以及办公成本开支难以维继，从而使公司陷于资金链断裂，公司面临破产的局面。

    现金流量管理的目标是否明确？

    要有效加强咨询公司现金流管理，为其经营发展创造良好条件，需从首先明确现金流量管理的目标开始。

    咨询公司主动进行现金流管理的目标包括：

    1保证公司日常运营合理开支；2增强现金流量的可预见性；3为不可预期的现金支出提供合理的内外部融资支持；4为闲置资金进行有效增值；5为公司业务拓展提供决策支持；6信用管理与融资能力管理。

    建立一个有效的财务控制体系

    建立一个有效的财务控制体系也是加强现金流管理的一个必要工具，可以从三方面来进行控制，让你的现金流管理拉响警报。

    ●人力成本控制

    人力成本控制应注重对人力资源的有效激励和成本支出的相对可控。人力成本主要包括工资、奖金、项目分成、福利费用、日常费用开销、特殊费用开销等。其中，对于工资、奖金、福利费用、日常费用开销等，属于可预期的日常人工成本，应在公司现金流安排中优先保证，对于奖金、分成、特殊开销等不属日常人工成本，存在一定程度的预期偏差，但公司应根据有关员工提出的建议及实际相关项目时间安排节点，纳入年度或月度计划，从计划上给予事先保证。

    ●项目成本控制

    项目成本控制是咨询公司实际赢利和现金收支平衡的重要手段。项目成本控制不仅责任在于公司财务部门，更在于培养业务部门在项目运作中的收入成本意识，既要有效节约成本，又要切实保证项目成本合理开支所需现金。

    ●现金预算制度

    咨询公司同样应加强现金流管理的计划性，为此可由财务部门牵头，各业务部门参与，以人和项目为基本要素，建立现金预算制度。由各业务部门编制现金收支计划，由财务部门汇总并监督执行。现金收支计划包括年度计划与月度计划。年度计划主要为公司整体业务进行提供决策参考，月度计划主要为项目实际进程中的有关收支提供具体支持。

    为经营发展创造良好条件

    健康的现金流量管理可以从经营、投资、融资等三方面为业务发展提供支持：

    ●对经营活动的支持

    1公司经营活动提供基本的现金流量保证；2公司进行业务拓展和市场扩张提供资金决策依据。

    ●对投资能力的支持

    1通过灵活运用各种短期投资工具，既保证资金的流动性，又使闲置资金有效增值；2通过对公司积累资金的充分利用和有关项目或操作平台的财务分析，寻求长期投资机会，帮助公司实际业务经营的多元化。

    ●对融资能力的支持

    1通过现金流的良好管理，实现公司财务指标的优化，提升企业财务资信；2通过积极拓展内源融资以外的其他外部融资，提升公司资金调度能力，合理利用财务杠杆，最大化公司价值。
 % 许多咨询公司认为
    
    资金管理（Funds Management）的原则主要是：划清固定资金、流动资金、专项资金的使用界限，一般不能相互流用；实行计划管理，对各项资金的使用，既要适应国家计划任务的要求，又要按照企业的经营决策有效地利用资金；统一集中与分口、分级管理相结合，建立使用资金的责任制，促使企业内部各单位合理、节约地使用资金；专业管理与群众管理相结合，财务会计部门与使用资金的有关部门分工协作，共同管好用好资金。

\subsection {资金管理中的常见问题}

    \subsubsection {资金管理意识淡薄，管理体系不健全}

    一些企业领导资金管理意识淡薄和没有资金时间价值观念。企业有钱时，不知如何规划使用，没有钱时就发愁，没有一种长期预算资金管理意识。在资金的整个循环过程中缺乏科学性和统一协调性，财务管理工作中没有充分考虑时间价值问题。现金流量是企业筹资、用资的关键。资金管理主要包括:余额管理和资金流量管理。一些企业对资金管理认识不到位，片面性追求产量和产值。尤其是在当前金融危机背景下，企业应对产品开发和未来风险进行合理评估和日常监控，避免产品占用大量资金和投资于风险大的项目。

    \subsubsection {资金管理模式不适应企业实际}

    有些企业没有建立完善的资金管理模式，有的建立了，但是不适应企业的实际。一般来说，资金管理模式分为两种：资金集中管理模式和分散模式。而企业一般分为单一企业和集团式企业。企业应根据实际情况建立适合企业资金管理的资金管理模式。一般来说集中式的资金管理对资金的集中控制和统一调配有利，但不利于发挥成员企业或分公司的积极性。成员企业或分公司在资金上过分依赖企业集团，若配套措施不到位，可能影响资金的周转速度，影响其对市场的应变能力。分散式的资金管理，有利于调动成员企业或分公司的积极性，但又难以避免资金分散、资金使用率低、沉淀资金比例大、资金使用成本高等固有缺点。由于集中式和分散式固有的不足，企业应有着重，结合两者模式的优点，根据企业的实际动态的混合的选用。我国企业的资金管理模式不科学，如有的集团公司内部还没有成立资金中心，没有运行资金集中管理机制；子公司多数没有大局意识和全盘观念，反映在资金日常管理上是五花八门、各自为政、存贷双高，出现了一些真空地带。

    \subsubsection {资金管理手段落后，使用效益低下}

    随着各个企业的迅速发展，跨区域经营的企业越来越多，这就出现如何对企业资金进行控制管理的问题。有的企业由于资金管理手段落后，资金控制能力不足，使得资金使用效益低下，在企业经营过程中付出了沉重的代价。尤其是大型集团公司由于成员企业众多，地域分布广泛，资金管理失控，监控缺乏手段，资金使用率低。其突出表现是:由于企业集团缺少统一资金管理系统，各个子公司、各种业务对资金流动的影响没有形成相关联的完整信息，难于有效监督，风险较大；由于缺少资金统一管理和规划使用，有的企业资金出现缺口，只能向银行贷款，而有的企业资金富裕，却无法给缺少资金的企业使用，只能向银行贷款，增加企业的财务费用，这导致了整个企业集团的资金成本上升，资源浪费和资金使用效益低下。这些问题严重制约了企业的健康发展，在当前这种形势下，甚至有可能威胁到企业的生存。

    \subsubsection {资金风险管理不足，引发严重财务风险}

    当前形势下，企业系统和非系统风险较大。而有些企业对资金风险管理没有引起足够的重视，主要表现在财务制度不健全，账户管理混乱，资金调拨不按流程和权限办理，造成企业资金的损失；违规对外担保或连环担保，给企业带来不必要的经济责任和法律责任，从而导致企业深陷担保诉讼泥潭。

\subsection {金融危机的影响}

    \subsubsection {市场需求萎缩，企业库存增加}

    金融危机引发的全球经济增速放缓，投资者信心不足，一方面造成企业投资不足，生产资料需求降低，企业效益下降，另一方面是消费者储蓄动机加大，减少了个人消费。这造成了整个市场的市场需求大幅萎缩，企业销售不畅，企业库存积压，同时由于市场价格加速下行，导致企业经营风险陡增，资金周转速度缓慢，资金链不能有效运行，也从不同程度上影响了企业的经营业绩。

    \subsubsection {企业间接融资难}

    根据《证券市场周刊》测算，在本次金融危机中我国有些金融机构海外投资出现了亏损，其中中行亏损额最大，约为38.5亿元。建行、工行、交行、招行及中信银行依次亏损5.76亿元、1.20亿元、2.52亿元、1.03亿元、0.19亿元。这些巨额的损失使各大商业银行在今后的发展中提高了自身经营的谨慎性，一方面严格规范和执行借贷标准，加强内部控制制度建设，提高自身管理能力；另一方面，虽然2009年我国实施积极宽松的货币政策，增加信贷规模和货币投放量，同时根据近日中国人民银行公布数据显示，新增贷款逐步回升，但是受金融危机持续升温、企业市场化经营形势恶化的影响，银行惜贷现象依旧，短期内难以见效，企业融资前景不容乐观，金融危机环境下企业资金紧张的矛盾还将继续。

    \subsubsection {企业通过直接融资获得资金更加困难}

    受金融危机的影响，我国股票市场受到冲击，导致了上证指数从2007年10月份的6124点暴跌至1638点，虽然2009年上证指数在回升，但是一直在3000点左右徘徊。金融危机打击了投资者的信心并在金融市场产生了示范效应。这使我国企业发行股票或者债券进行融资更加困难，企业通过直接融资缓解企业营运资金的渠道受阻。

    \subsubsection {金融危机引发了企业资金链断裂危机}

    所谓资金链，是指维系企业正常生产经营运转所需要的基本循环资金链条。资金链犹如一个企业的血脉，其循环过程可以简短的表述为：现金——资产——现金。企业要生存发展，就必须保持这个循环良性运转。当企业在资产难以变现或根本没有资产可以变现的情况下，由于现金流通不畅或没有可支配现金不能偿付到期债务，由此引发资金链断裂的危机。由于当前企业出现了直接融资、间接融资难，产品需求下降以及企业内部经营管理方面的原因造成企业资金链断裂。

\subsection {参考解决措施}

    \subsubsection {增强资金管理意识}

    首先，企业应树立资金进行统一管理的观念，不管是单一企业还是集团企业，在当前形势下，要对资金进行统一管理，规划使用。其次，树立现金流量的观念，在财务管理的具体工作中，为管理人员提供现金流量的信息，除年终提供的现金流量表之外，在日常工作中可根据不同情况，编制现金流且计划，以及短期现金流量预测报告和长期现金流量报告。要加强现金流量的控制和分析，严格控制现金流入和流出，保证企业始终具备支付能力和偿债能力。最后，提高企业资金管理的风险意识，要充分估计各个项目风险，谨慎投入资金。

    \subsubsection {全面提升企业资金管理水平}

    在当前形势下，不管是单一企业还是集团企业，资金集中管理是发展的必然趋势。实施企业资金集中管理，对于企业的生存与发展具有重要作用。它有助于企业完善整体资金链，实现整体利益的最大化，有利于集中进行战略方向的调整，有效地降低企业控制成本，提高企业整体信用等级，降低财务成本，优化资金管理体系，提高资金的使用效率。目前，一旦企业进行资金集中管理，企业可以按照自身需求对各个子账户的资金进行相应地归集，加强内部资金的整合和统筹管理，实现内部资金的相互平衡，提高资金的使用效益。

    \subsubsection {积极利用计算机技术，保证资金管理}

    随着计算机信息技术的发展，企业在实施资金的集中管理和监控中可以大力应用计算机技术。在企业中大力利用新的资金管理手段和模式，而计算机网络技术和统一的财务管理软件是先进的管理思想、管理模式和管理方法的有效载体，也是实施资金集中管理和有效监督控制的必然选择。其次，借助ERP系统的优势，提升企业资金管理效率。ERP 系统简化了采购、销售与财务之间的流程过程，充分利用财务与其他业务之间数据信息的互通，提高管理效率，使资金管理贯穿于整个企业业务流程的每个环节，对企业各个环节进行实时监控，有效发挥财务监管机制。

    \subsubsection {强化财务监督}

    首先，加强内部管理，对合同事前审批、事中执行、事后评价，进行严格地监控与分析。增强风险防范意识，建立风险预警机制。对经营业务的采购、销售、库存等各风险节点资金使用情况进行认真梳理，实时监控可能发生的风险，有效应对，将风险控制在最低。其次，积极开展内部审计，前移监督关口。企业的内部审计是严格监督、考核企业财务资金管理的重要环节，是强化监督约束机制、使预算取得实效的保障。健全内部审计监督考核制度，保证企业财务信息的真实可靠。变过去的“事后监督”为事前、事中监督和适时监督，围绕企业的发展目标和年度预算，对公司资金流向、财务状况变动等情况实施全过程的跟踪和监控，定期检查，及时反馈，防微杜渐，确保预算的严肃性和企业发展目标的如期实现。通过以上两方面的努力保证企业资金安全和完整。

    \subsubsection {拓宽融资渠道}

    有条件的企业积极取得政府和金融机构的支持，获得资金支持。同时企业应根据自身的实际情况，积极拓展融资渠道。在内源融资方面可以实施员工持股计划、预收账款融资、质押应收账款融资、专利权质押融资以及产权交易等内源性融资方式。在外源融资方面，除了依赖传统外部资金支持，如银行存款、风险基金、发行债券或企业上市等，努力寻求企业间金融互助合作，有条件的企业可以考虑风险投资、融资租赁和股权融资等创新型的融资方式 。
 % 资金管理中的常见问题
    \subsection {浅析“以财务为中心”}

    不少公司都谈到类似“企业管理以财务为中心，管理信息化应当先上财务系统”的论点。抱有这类论点的个人和公司很多，其中不乏由电算化会计转型到 ERP 系统的公司，甚至有些外国公司的发言也是如此讲。

    企业管理涉及的业务很多，其中有一些属于核心业务，如，制造业的产品研发、销售、生产、采购、财务等; 没有这些业务就谈不上企业运作。 就 ERP 内部集成而言就包括了所有的核心业务。财务与其他业务一样，都属于核心业务。任何类型的企业，核心业务的内容可能会有变化，但却都少不了财务; 也许，这正是人们容易把财务作为“中心”的一个原因。

    但是我们要注意到，每一项核心业务都是一种流程，而不同流程的发生是有先有先后的，就是项目管理中经常强调的“先导”与“后续”的关系。在制造业，财务成本又是销售、生产和采购3项核心业务运作的结果，相对于销售、生产和采购来讲是一种“后续”流程。通过财务分析再反过来，指导或修正经营业务。 这里的信息流程是闭环和“互动”的，可见，各个业务系统之间是一种“相互依存、相互作用”的集成关系，而不是以哪个为中心的关系。

    财务管理大量的工作是账务处理，最后生成国家财政部要求的3种报表——损益表、资产负债表和现金流量表。这些报表反映的是一个时期企业经营状况的结果，是在一定的会计期末向企业外部公开发布的信息。这里的关键词，一个是“结果”，一个是“期末”，既不能反映 “过程”，又不能做到 “实时”。

    在其他业务系统还没有实施之前，所有进入财务系统的原始数据，都是手工录入的，甚至是“人工合成” 的，不能体现信息集成，实时性和准确性都得不到保证。报表上所有的数据都要来自各个业务系统，没有供销产等各种业务信息，财务系统将成为无源之水，无本之木。

    业界有的人士说得对 “财务报表只是一个结果，只知道结果而不了解过程，对制造业改进管理来说没有多大用处的”。只看结果，不看过程，是无法控制和优化流程的。与实现信息化的初衷相悖，这本是很容易理解的浅近道理，但目前社会上的认识却非常混乱。

    因此，以“财务为中心”的提法会影响企业依据自身的特点实事求是地确定实施 ERP 系统的步骤。而且，如果作为软件设计的指导思想，会成为束缚产品的发展的思想桎梏。

    不是说企业在实施 ERP 系统的进程中，不可以先上财务系统，不能绝对化。但是，如果先上财务系统，就必须对上面提到的现实情况要有一个清楚的认识。
 % 浅析“以财务为中心"
    \subsection {精细化资金管理}

    \subsubsection {为什么要对资金进行精细化管理}

    很简单，因为资金的节省就是净利润。因为省钱比挣钱要容易得多。对于企业来说你实现100元的销售收入可能只获得10元的利润，但如果你省10元，10元同样也是利润，但付出的努力是完全不一样的。

    另外对资金进行精细化管理的另一个原因是理解资金和业绩的关系。比如销售人员的资金和销售业绩的关系，研发资金和技术成果的关系等等。通过对这种关系的分析从而更好的配置有限的资金预算。

    精细化资金可以进行业绩评估。尤其同等资金的产出指标。

    \subsubsection {目前资金管理的模式}

    目前可以说所有的公司都是进行大呼隆式的的粗放式管理，包括年度的资金预算和实际资金的核算。顶多核算到部门就算了，资金成了流水账。很多的信息都被淹没了。

    我们可以参考SAP的资金核算大体如下：

    年度预算：可以针对部门按照成本中心和成本要素进行预算，另外也可以使用内部订单进行一些项目预算。

    日常核算：所有的资金都是记到成本中心或内部订单，简单的流水记账，如果你想知道某种资金或个人的资金就比较麻烦。可能需要导出数据做很多处理。

    可以说从预算和实际核算都是粗放型的。为什么，没有按照部门，人员，业务活动和业绩进行。最终也是糊涂账。

    \subsubsection {精细化的资金管理模式}

    对于一些资金处理较多的组织和部门，比如销售公司或者每月需要处理很多资金的财务中心，同时需要对不同的资金进行分析和评估，楼主推荐如下模式：

    \begin{enumerate}
        \item  报销人员网上申报个人资金（包括年初的预算和实际的报销），主管人员审批预算和实际报销并评估业务，月度资金报告和业绩报告。资金系统用工作流驱动整个流程，可以结合移动商务时时批准，减少高层人员出差的审批问题。同时实际的资金单据全部使用条形码管理，十分便于财务人员集中处理。

        \item  数据接口方式。资金系统实现与各种现有ERP系统的无缝转接。采用批处理方式，资金系统可以实现输出用户需要的格式进行数据导入，方便的实现数据录入，没有任何冗余动作和风险。

        \item  简单明了的驾驶舱分析报告。高层管理人员可以方便的调阅部门和人员的资金执行报告，随时进行管理和调控。
    \end{enumerate}

    总之，精细化的资金管理就是要从目前粗放型的管理实现由部门和项目的管理更象预算，个人和资金类型的多纬度的精细管理演进。
 % 精细化资金管理
    \subsection { 营运资金优化}

    四种加强企业营运资金管理方式：

    一、通过对购销环节中的风险进行事先预测、分析和控制，尽量规避风险

    许多企业为了实现利润、销售更多产品，经常采用赊销形式。片面追求销售业绩，可能会忽视对应收账款的管理或管理效率低下。例如对赊销对象的现金流动情况及信用状况缺乏，以及未能及时催收货款，容易出现货款被客户拖欠，造成账面利润很高，但实际资金却很少。对此，财务部门应加强对赊销和预购业务的控制，制定相应应收账、预付货款控制制度，加强对应收账款的管理，及时收回应收账款，减少风险，从而提高企业资金管理使用率。

    二、控制成本费用，提高企业利润，增加企业价值

    会计利润是当期收入和费用成本配比的结果。在任何收入水平下，企业都要做好对内部成本、费用的控制，做好预算，加强管理力度，减少不必要的支出，才能够提高利润，增加企业价值。

    三、加强企业财务预算，提高企业营运效率

    财务管理应站在企业全局的角度，构建科学的预测体系，进行科学预算。预算包括销售预算、采购预算、投资预算、人工预算、费用预算等等，这些预算使企业能预测风险，及时得到资金管理的各种信息，及时采取一定措施，防范风险，提高效益。同时，制定的各种预算，不仅可以协调企业各部门的工作，避免部门间冲突，提高内部协作的效率。而且，销售部门在销售、费用等预算指导下，还可事先对市场有一定了解，把握市场变化，减少存货的市场风险。

    四、完善财务制度建设，进一步加强营运资金管理

    \begin{enumerate.zh}
        \item  明确内部管理责任制。

        很多企业认为催收货款是财务部门的事，与销售部门无关，其实这是一种错误的观点。事实上，销售人员对催收应收账款应负主要责任。如果销售人员在提供赊销商品时，还要负收回应收账款的责任，那么，他会谨慎对待每一项应收账款。销售人员应在征得财务部门的认可后才能发货，并且要按购销合同及时收回货款。

        \item  建立客户信用档案。

        企业应在财务部门中设置风险控制员，应通过风险控制员对供应商、客户的信用情况进行深入调查和建档，并进行信用等级设置，对处于不同等级的客户实行不同的信用政策，减少购货和赊销风险。风险管理员对客户可从以下方面进行信用等级评定：考察企业的注册资本;偿还账款的信用情况;有没拖欠税款而被罚款的记录;有没有拖欠供货企业货款的情况;其他企业的综合评价。风险管理员根据考察结果向总经理汇报情况，再由风险管理员、财务部门经理、销售部门经理、总经理讨论后确定给予各供应商及客户的货款信用数量。如果提供超过核定的信用数量时，销售人员必须取得财务经理、风险管理员及总经理的特别批准。如果无法取得批准，销售人员只能降低信用规模或者放弃这桩业务，这样就能控制销售中出现过度多的应收账款现象，减少风险。

        \item  严格控制信用期。

        应规定应收账款的收款时间，并将这些信用条款写进合同，以合同形式约束对方。如果出现未能在规定时间内收回应收账款，企业可依据合同，对拖欠货款企业采取一定措施，以及时收回货款。

        \item  通过信用折扣鼓励欠款企业在规定时间内偿还账款。

        很多企业之所以不能及时归还欠款是因为他们及时归还得不到什么好处，拖欠也不会有什么影响。这种状况会导致企业应收账款回收的效率比较低下。为了改善这种局面，企业可以采取相应的鼓励措施，对积极回款的企业给予一定的信用折扣。

        \item  实施审批制度。

        对不同信用规模、信用对象实施不同的审批级别。一般可设置三级审批制度。由销售经理、财务经理和风险管理员、总经理三级审核。销售部门如采用赊销方式时，应先由财务部门根据赊销带来的经济利益与产生的成本风险进行衡量，可行时再交总经理审核。这样可以提高决策的效率，降低企业经营的风险。

        \item  加强补救措施。

        一旦发生货款拖欠现象，财务部门应要求销售人员加紧催收货款，同时风险管理员要降低该企业的信用等级;拖欠严重的，销售部门应责令销售人员与该企业取消购销业务。

        \item  建立企业内部控制制度。

        主要包括存货、应收账款、现金、固定资产、管理费用等一系列的控制制度。对违反控制制度的，要给予相关责任人一定的惩罚。

        \item  严格控制开支，对各种开支采用计划成本预算。

        对各种容易产生浪费的开支要采取严格的控制措施。例如，很多企业业务招待费在管理费用中占很大比例，在计征所得税时无法全额在税前扣除。对此，企业应该要求销售人员控制招待费支出，由财务部门按其月销售收入核定适当的招待费标准。
    \end{enumerate.zh}

    总之，营运资金管理在企业销售及采购业务中处于重要地位，对企业利润目标的实现会产生重大影响。营运资金管理应是对销售工作的控制而不是限制，促进销售部门减少销售风险，提高利润水平。企业领导人不仅应重视财务管理，还要协同搞好资金营运管理。
 % 营运资金优化
    \subsection {资金风险管理}

    资金风险的管理技巧，是投资成败的关键之一。很多非常聪明的投资者，都是在资金管理上犯了情绪化的错误。我们建议您：
    \begin{enumerate}
        \item  在资金管理上最重要是保障本金的安全，其次才是获利最大化的追求。
        \item  止损/止盈纪律的执行要坚决实行程式化。
        \item  失误后要先执行纪律，再进行场外反思。
    \end{enumerate}

\subsubsection {根据自身情况调节风险}

    如果您是刚开始交易的新手，希望您掌握好账户资金以及开仓所占保证金的比例。

    一般风险规则告诉我们：每笔交易的风险不能超过5\%-15\%。比如说，你用所有保证金的5\%开仓位。如果损失，你就用剩下的95\%的5\%来开仓位。

\subsubsection {风险回报率}

    目标盈利和所能承受损失比率一般设为2：1，但3：1更好。

\subsubsection {止损}

    不要忘记止损。通常，当市场跟您预期的背道而驰时，投资者总是不自觉地想调整止损位。但是止损一旦设定就请不要更改，否则止损就失去了保护资产的作用。

    我们鼓励每个投资者在场外反思如下问题：

    \begin{enumerate}
        \item  我能承受多大风险？
        \item  我最多要赚多少，最多能亏多少？
        \item  目前市场情况如何，很动荡还是很平静？
        \item  我做这个仓位的逻辑是什么？
        \item  需要多长时间可以下结论我做一个仓位的逻辑是否正确？
    \end{enumerate}
 % 资金风险管理
    融资租赁

    随着国家对融资租赁行业的支持力度加大，融资租赁企业如雨后春笋般涌现，发展迅速的同时，面对的国内外竞争对手也越来越多，业务操作和管理的复杂性加大，强化行业的管理能力势在必行。近年来，科学管理受到普遍重视，特别是我国加入WTO后，受到全球经济一体化的挑战，越来越多的企业认识到信息化建设的迫切性。信息化对提升企业管理水平的重要性已经经国际众多大型企业验证，融资租赁在我国作为一个新兴行业，要想进一步发展和提高管理能力，就必须加强行业的信息化建设。

     加快信息流转，突破资源壁垒

    利用信息系统将客户信息、业务信息、设备信息、合同信息、资金信息等项目信息统一管理，形成一个完整的企业信息资源库。当项目开始启动，整个信息资源库就会被调用起来，消除了以前人工管理引发的信息不通畅，资源无法充分利用的被动局面，信息流动始终贯穿企业经营活动的整个过程。

    顺畅业务流程，提高工作效率

    融资租赁业务涉及多个业务对象，如租赁公司、设备供应商、客户、代理商等，每个环节的业务流程都有其运转的规律，相互之间缺乏对接。利用IT手段将专业的融资租赁系统与企业内部系统，如1104系统、CRM系统、财务系统、网银系统等整合起来，可简化和优化企业的业务流程，某个系统中的流程结束后，将自动触发另一个流程，形成闭环顺畅的业务处理程序，提高工作效率，增强企业的业务处理能力。

    规范业务处理，控制项目风险

    融资租赁的业务处理，对内要符合企业的管理理念和业务制度，控制项目风险；对外要遵循国家监管机构（如银监会等）制定的各类法律规范，降低行业风险。是以，企业需要一套深刻融合内外部的各类规章制度的业务管理系统，通过表单、流程、模块等将其规范到日常的工作中，通过与监管部门的系统实现对接，来规范业务处理。系统需要具备随需应变的功能，如自定义流程、灵活的权限管理等来应对与时俱进的法律法规及企业的业务规范，使项目始终业务合规、风险可控。

     数据报表分析，提供决策支持

    由于融资租赁行业的特殊性，在项目运作过程中，需要丰富多样的业务报表来对管理层的决策提供支持，如行业逾期率图表、时序曲线图、逾期租金清单、授信额度使用率报表、租息差额汇总等等，管理层需要通过这些报表上的信息了解项目进展，并制定下一步的决策。缺乏信息化管理的企业，无法实时获取数据报表，且数据报表的结论具有人为的主观性。一套了解融资租赁行业业务规则和风险机制的业务管理系统能为企业管理决策提供实时、客观、科学的数据报表，使企业的决策高效高质。

    多语言多架构，助力国际化发展

    随着融资租赁行业进入快速发展期，外资进驻中国市场，内资投资国外市场，融资租赁行业的国际化趋势越来越明显，企业的员工组成和组织架构的复杂程度也随之增加。企业的信息化系统需要能支持不同的语言环境，如简体中文、繁体中文、英文、日文等，以此满足不同地区员工的工作习惯。系统中的权限管理要能与企业的组织架构相匹配，以此来实现横向纵向的矩阵式管理，既满足管理需求，也满足业务需求，助力企业国际化发展。

    通过建设信息化系统促进融资租赁企业的健康生存和发展，提高企业的管理能力和业务水平，是企业发展的必由之路。信息化能帮助融资租赁企业实现资源整合，流程严谨，业务合规，项目监控，管理灵活，为融资租赁企业提供一个风险防控平台，协助融资租赁企业在业务效益和管理效益上收获双赢！
 % 融资租赁
    \subsection {资金管理的误区}

    “现金为王”这样的经营理念慢慢的被更多的企业家所接受,大家对如何做好资金管理也都有自已的看法.在财务管理的理论中,也有很多计算模型和观念.其中有一点说企业最佳的现金持有状态是_通过计算企业最小的现金持有量,除此之外,账上不保留更多的钱,但是需要的时候就可能拿到钱.

    理论有的时候只是计算一个理想化状态,很多时候它只存在于书本和考试中.可是有很多人迷信书本,在现实的企业管理追求出现这种理想状态,认为只有这样才是企业资金管理的最佳状态,这是非常可怕的.

    我们都知道,企业自身永远是处于一个动态发展的过程中,同时,企业也处于社会动态变化的过程中,在二个动态过程里面,找到一个静态或不断变化的静态理想状态,几乎是不可能的事情.所以在实践中,我们可以看到,那些活得最久的企业,它们的资金管理对比理论来说,都不是最佳的状态.华人首富李嘉诚企业永远保留大量的现金,远大在很多年前就宣称企业有5亿现金是不动的.

    对远大这种做法的质疑有很多,认为企业这种做法失去了很多发展的机会.这种说法当然有它的道理,包括现金是收益性最差的资产,这种说法都是对的.

    但是,在说这些话的时候,很多人忘记了一点:任何事情都有它的二面性.保有大量现金,失去发展机会和收益差的同时,这也是企业长久发展最安全的方式,企业真的很难做到二者的兼顾.

    在这个时候,往往是考验企业家智慧的时候,你把企业做成了短跑英雄,还是做一个长跑的冠军.我们眼前闪过了太多的短跑英雄,顺驰、飞龙、三株等等。

    在一些企业管理的核心理念上具有什么样的态度,决定了企业发展的心态.做为企业管理核心的资金管理上,更是如此.企业其实就是一项马拉松,短跑心态可能有开始的风光,但那一定是流星一样的短暂,笑到最后才是最好的!
 % 资金管理的误区

\practices
\subsection{会计科目管理}
管理现有的会计科目
\opset{会计科目管理}{
    \item \ops{新增会计科目}{
        \item  点击\button{新建}按钮，进入编辑会计科目界面
            \screenshot{1.png}
        \item  输入\button{编号}及\button{科目名称}
        \item  选择\button{上级科目}（图）
        \item  点击\button{保存}按钮，保存新的会计科目
    }
    \item \ops{查看会计科目}{
        \item  选择需要查看的会计科目
        \item  点击工具条上的\button{查看}按钮
        \item  查看会计科目的详细内容
    }
    \item \ops{修改会计科目}{
        \item  选择需要修改的会计科目
        \item  点击工具条上的\button{修改}按钮，进入编辑会计科目界面
        \item  修改需要修改的会计科目明细
        \item  点击\button{保存}按钮，保存更改后的会计科目
    }
    \item \ops{删除会计科目}{
        \item  选择需要删除的会计科目
        \item  点击工具条上的\button{删除}按钮
        \item  在弹出的对话框中点击确定按钮，删除对应会计科目
            \screenshot{2.png}
    }
}

\subsection{资产初始化}
初始化资产管理,设置企业出世资产
\opset{资产初始化}{
    \item \ops{新建初始资产单}{
        \item  点击\button{新建}按钮，进入新建初始资产单界面
            \screenshot{9.png}
        \item  点击\button{添加}按钮，弹出添加初始资产单项目对话框
        \item  选择\button{会计项目}
        \item  选择\button{二级会计项目}
        \item  输入\button{金额}
        \item  选择\button{币种}
        \item  点击\button{确定}按钮，添加初始化资产单项目
        \item  点击\button{保存}按钮，保存初始化资产单
    }
    \item  \ops{修改初始资产单}{
        \item  选择需要修改的初始化资产单
        \item  点击工具栏上的\button{修改}按钮
        \item  修改初始化资产单的相应内容
        \item  点击\button{保存}按钮，保存更改后的初始化资产单
    }
    \item \ops{删除初始资产单}{
        \item  选择需要删除的初始化资产单
        \item  点击工具栏上的\button{删除}按钮，在弹出的对话框中点击确定，删除对应的初始化资产单
            \screenshot{10.png}
    }
}

\subsection{资金流业务管理}
\opset{资金业务管理}{
    \item \ops{新建资金流业务单}{
        \item  点击工具栏上的\button{新建}按钮，进入新建资金流业务单界面
            \screenshot{6.png}
        \item  输入\button{金额}，\button{摘要}，\button{业务详细}
        \item  点击\button{保存}按钮，保存新的资金流业务单
    }
    \item \ops{修改资金流业务单}{
        \item  选择需要修改的资金流业务单
        \item  点击工具栏上的\button{修改}按钮，进入到修改资金流业务单界面
        \item  修改相关内容
        \item  点击\button{保存}按钮，保存更改后的的资金流业务单
    }
    \item \ops{删除资金流业务单}{
        \item  选择需要删除的资金流业务但
        \item  点击工具栏上的\button{删除}按钮
        \item  在弹出的对话框中点击\button{确定}按钮，删除对应资金流业务单
            \screenshot{7.png}
    }
}

\subsection{凭证管理}
凭证管理用来管理资金流业务单，实现收支平衡
\opset{凭证管理}{
    \item \ops{新建凭证}{
        \item  点击\button{新建}按钮进入新建凭证界面
            \screenshot{3.png}
        \item  选择对应的资金流业务单
        \item  点击\button{添加}按钮添加凭证项目，确保借贷平衡
            \screenshot{4.png}
        \item  点击\button{保存}按钮，保存新凭证
    }
    \item \ops{修改凭证}{
        \item  选择需要修改的凭证
        \item  单击\button{修改}按钮
        \item  修改需要修改的项目或对应的资金流业务单
        \item  点击\button{保存}按钮，保存修改后的凭证
    }
    \item \ops{删除凭证}{
        \item  选择需要删除的凭证
        \item  点击\button{删除}按钮
        \item  在弹出的对话框中点击\button{确定}按钮，删除凭证
            \screenshot{5.png}
    }
}

\subsection{资产查询}
查询当前资产
\screenshot{8.png}
\opset{资产查询}{
        \item \ops{查询步骤}{
            \item  勾选全部选择框，点击\button{查询}按钮，查询全部资产
            \item  没勾选全部选择框的情况下
            \ul {
            \item  点击按钮选择会计科目
            \item  点击\button{查询}按钮，查询相应会计科目下的资产情况
            }
    }
}
